\documentclass[conference]{IEEEtran}
\usepackage{gvv-book}
\usepackage{gvv}
\usepackage{algorithm}
\usepackage{algpseudocode}
%\usepackage{listings}
\usepackage{xcolor}
%\usepackage{amsmath}
\usepackage{graphicx}
\usepackage{float}
\usepackage{booktabs}
\usepackage{url}
\usepackage{circuitikz}
\usepackage{hyperref}
\usepackage{tabularx}
\lstset{
  basicstyle=\ttfamily\small,
  keywordstyle=\color{blue}\bfseries,
  commentstyle=\color{green!40!black},
  stringstyle=\color{red},
%  numbers=left,                     % keep numbers
%  numberstyle=\scriptsize\ttfamily, % more readable numbers
%  numbersep=8pt,                    % add spacing from code
  xleftmargin=1.5em,                % indent code so numbers don't overlap
  frame=single,
  breaklines=true,
  breakatwhitespace=true,
  tabsize=2,
  showstringspaces=false,
  captionpos=b,
  keepspaces=true,                  % preserve spacing properly
  linewidth=\columnwidth,           % constrain to column width
  columns=fullflexible              % prevent box overflow
}
\def\UrlBreaks{\do\/\do-\do\_\do\.}
\title{ESP-Flasher}
\author{
	\IEEEauthorblockN{Darsh Gajare-EE25BTECH11020 and Nishid Khandagre -EE25BTECH11043}
}
\begin{document}
\maketitle

\tableofcontents

% ─────────────────────────────────────────────────────────────
\section{Uploading Code to Arduino Using ESP32 Wirelessly}

\subsection{Introduction}

Uploading code to an Arduino through a phone can be
done using applications like ArduinoDroid; however uploading
solely through Termux is different because Termux does not
have permission to directly access the USB port on Android
devices.

This document describes a hardware workaround: an ESP32
acts as a transparent bridge between WiFi and UART, forwarding the STK500 programming protocol from the phone to the
Arduino bootloader wirelessly.

\subsection{Overview}

The phone compiles the sketch using PlatformIO and invokes avrdude, which opens a TCP socket to the ESP32 on
port 328. The ESP32 forwards all bytes transparently between
the socket and UART2, relaying the STK500 protocol to the
Arduino bootloader.

\begin{table}[H]
\centering
\resizebox{\columnwidth}{!}{%
\begin{tabular}{ll}
\toprule
Component & Role \\
\midrule
Phone (Termux + Debian)  & Compile via PlatformIO; upload via avrdude \\
ESP32 NodeMCU            & WiFiServer TCP:328; UART2 relay at 115200 baud \\
Arduino UNO (ATmega328P) & Optiboot bootloader at 115200 baud \\
\bottomrule
\end{tabular}%
}
\end{table}

\subsection{Hardware Connections}

\begin{table}[H]
\centering
\resizebox{\columnwidth}{!}{%
\begin{tabular}{llll}
\toprule
ESP32 Pin  & Arduino Pin & Dir. & Note \\
\midrule
5V         & VIN         & $\rightarrow$     & Power via onboard regulator \\
GND        & GND         & $\leftrightarrow$ & Common ground --- essential \\
P17 (TX2)  & Pin 0 (RX)  & $\rightarrow$     & Direct wire \\
P16 (RX2)  & Pin 1 (TX)  & $\leftarrow$      & Via 3$\times$10 k$\Omega$ voltage divider \\
P4         & RST         & $\rightarrow$     & Reset control \\
\bottomrule
\end{tabular}%
}
\end{table}

\subsubsection{Pin Wiring}

\subsubsection{Voltage Divider}

Arduino TX outputs 5 V logic; ESP32
GPIO pins support 3.3 V. A resistor voltage divider on the
P16 (RX2) line is required:

\begin{align}
  V_{out} = V_{in} \times \frac{R_2}{R_1 + R_2}
           = 5\,\text{V} \times \frac{10\,\text{k}}{30\,\text{k}}
           = 3.3\,\text{V}
\end{align}

\subsubsection{Reset Capacitor}

A 100\,$\mu$F electrolytic capacitor connected between Arduino RST (+) and GND ($-$) extends the
bootloader window to 1--3 seconds, giving time to press the
reset button on the Arduino.

\subsection{ESP32 Firmware Flashing Using ArduinoDroid}

The ESP32 firmware starts a TCP server on port 328 and
relays bytes between the WiFi socket and UART2. Code:
\begin{quote}
\href{https://github.com/gajaredarsh/esp-flasher/tree/main/esp_programmer/src}{\texttt{esp32\_programmer/src/main.cpp}}.\end{quote}
Replace \texttt{YOUR\_SSID} and \texttt{YOUR\_PASSWORD} with your network credentials.

\begin{enumerate}
  \item Compile using Platformio
  \item Connect the ESP32 to the phone via OTG cable.
  \item Open ArduinoDroid $\rightarrow$ paste firmware $\rightarrow$ select ESP32 Dev Module.
  \item Tap Upload (hold BOOT on the ESP32 if upload fails).
  \item Use Serial Monitor at 115200 baud to read the IP address.
\end{enumerate}

\subsection{Arduino Project Setup}

\subsubsection{Create Project}
Use the commands
\begin{lstlisting}[caption={Initialize Arduino project}]
mkdir arduino_code && cd arduino_code
pio init --board uno
nvim src/main.cpp # write your sketch
\end{lstlisting}

\subsubsection{PlatformIO Configuration}

Paste the contents of
\begin{quote}
\href{https://github.com/gajaredarsh/esp-flasher/blob/main/arduino-code/platformio.ini}{\texttt{arduino\_code/platformio.ini}}
\end{quote}
into your project's \texttt{platformio.ini}. Replace \texttt{ESP\_IP}
with the ESP32's IP address (e.g.\ 192.168.43.100).

The configuration uses \texttt{upload\_protocol = custom}
with avrdude connecting over \texttt{socket://ESP\_IP:328} at
115200 baud, targeting atmega328p with the stk500v1 programmer.

\subsection{Uploading Firmware Wirelessly}

\begin{lstlisting}[caption={Build and upload}]
cd arduino_code
pio run --target upload
\end{lstlisting}

Press RST on the Arduino immediately after running the command.

\subsection{Finding the ESP32 IP Address}

The ESP32 IP address changes when switching networks.

ArduinoDroid Serial Monitor --- connect via OTG, open
Serial Monitor at 115200 baud, press RST on ESP32.

\subsection{The Timing Issue}

Uploading relies on the Arduino auto-reset mechanism.
When a serial connection opens, the DTR line drops, pulsing
the RESET pin and forcing the Arduino into bootloader mode.
The bootloader listens for only $\approx$1 second before timing out.

Over TCP, avrdude cannot control DTR automatically. The
100\,$\mu$F capacitor (Section \ref{sec:resetcap}) solves this by extending the
window to 1--3 seconds.

% ─────────────────────────────────────────────────────────────
\section{ESP32-to-ESP32 OTA Firmware Updates}

\subsection{Introduction}

This chapter extends the wireless programming concept to
ESP32-to-ESP32 over-the-air (OTA) firmware updates. Instead
of a WiFi-to-UART bridge, ESP32 \#1 acts as a WiFi Access Point
and HTTP OTA controller. ESP32 \#2 connects to it and
updates its own firmware wirelessly --- no USB required after
initial setup.

Method: HTTP OTA over raw TCP is used instead of
Arduino OTA. Arduino OTA relies on mDNS discovery and
has no built-in integrity verification, making it unreliable for
ESP-to-ESP flashing. HTTP OTA uses direct IP addressing
with MD5 checksum verification.

\subsection{System Architecture}

\begin{table}[H]
\centering
\resizebox{\columnwidth}{!}{%
\begin{tabular}{llll}
\toprule
Device     & Role                  & IP           & Port \\
\midrule
ESP32 \#1  & Sender / Access Point & 192.168.4.1  & 80   \\
ESP32 \#2  & Receiver / OTA Target & 192.168.4.2  & 8080 \\
\bottomrule
\end{tabular}%
}
\end{table}

\subsubsection{Network Topology}

\subsubsection{Data Flow}

\begin{enumerate}
  \item Phone POSTs firmware.bin to ESP32 \#1 at \texttt{/upload}.
  \item ESP32 \#1 saves binary to LittleFS and computes MD5.
  \item Phone GETs \texttt{/trigger} on ESP32 \#1.
  \item ESP32 \#1 opens raw TCP connection to 192.168.4.2:8080.
  \item ESP32 \#1 sends HTTP POST with binary body and \texttt{x-md5} header.
  \item ESP32 \#2 streams bytes into Update library, verifies MD5.
  \item ESP32 \#2 replies HTTP 200 OK.
  \item ESP32 \#2 calls \texttt{Update.end(true)} $\rightarrow$ sets boot partition $\rightarrow$ \texttt{ESP.restart()}.
\end{enumerate}

\subsection{Initial Setup}

\subsubsection{Flash ESP32 \#2 (Receiver) via USB}

Compile the following code using Platformio
\begin{quote}
\href{https://github.com/gajaredarsh/esp-flasher/blob/main/esp_receiver/src/main.cpp}{\texttt{esp32\_receiver/src/main.cpp}}
\end{quote}
	Using ArduinoDroid flash via USB. This is the only time
USB is needed for ESP32 \#2. On boot it connects to ESP32-OTA-Net and starts the OTA server on port 8080.

\subsubsection{Flash ESP32 \#1 (Sender) via USB}

Compile the following using Platformio
\begin{quote}
\href{https://github.com/gajaredarsh/esp-flasher/blob/main/esp_sender/src/main.cpp}{\texttt{esp32\_sender/src/main.cpp}}
\end{quote}
	Using ArduinoDroid flash via USB. ESP32 \#1 boots as
WiFi AP with SSID ESP32-OTA-Net, password 12345678,
IP 192.168.4.1.

\subsubsection{Verify Both Boards}

Connect your phone to ESP32-OTA-Net, then:

\begin{lstlisting}[caption={Verify both boards}]
curl http://192.168.4.1/status # ESP32 #1
curl http://192.168.4.2:8080/status # ESP32 #2
\end{lstlisting}

Both should respond with \texttt{\{"status":"ready",...\}}.

\subsection{Repeatable Update Workflow}

\subsubsection{Step 1 --- Compile Firmware Binary}
Using
\begin{lstlisting}[caption={Compile with PlatformIO}]
pio run
# binary at: .pio/build/esp2/firmware.bin
\end{lstlisting}

\subsubsection{Step 2 --- Upload Binary to ESP32 \#1}
Using
\begin{lstlisting}[caption={Upload binary to ESP32 \#1}]
# PlatformIO binary
curl -F "firmware=@.pio/build/esp2/firmware.bin" \
  http://192.168.4.1/upload
# Verify it landed
curl http://192.168.4.1/status # firmware_loaded: true
\end{lstlisting}

\subsubsection{Step 3 --- Trigger OTA Push}
Using
\begin{lstlisting}[caption={Trigger OTA push}]
curl http://192.168.4.1/trigger
\end{lstlisting}

ESP32 \#1 connects to 192.168.4.2:8080, sends the binary
with the MD5 header, and reports success. ESP32 \#2 verifies,
writes, and reboots.

\subsubsection{Step 4 --- Verify}

Wait $\approx$5 seconds for ESP32 \#2 to reboot, then:

\begin{lstlisting}[caption={Verify ESP32 \#2 is back online}]
curl http://192.168.4.2:8080/status
\end{lstlisting}

\subsection{Golden Rule}

Every sketch deployed to ESP32 \#2 must include
\texttt{handleOTAClient()} and \texttt{otaServer.available()}
in \texttt{loop()}. If you flash a plain sketch without it, ESP32 \#2
loses wireless update capability and requires USB to recover.
See the receiver template for the minimum required structure.

\begin{lstlisting}[caption={Minimum required structure for every ESP32 \#2 firmware}]
#include <WiFi.h>
#include <Update.h>

WiFiServer otaServer(8080);

// paste handleOTAClient() from esp32_receiver_final.ino

void setup() {
  connectWiFi(); // connect to ESP32-OTA-Net
  otaServer.begin(); // always start OTA server
  // your app setup ...
}

void loop() {
  WiFiClient c = otaServer.available();
  if (c) handleOTAClient(c); // always poll for OTA
  // your app code below ...
}
\end{lstlisting}


\end{document}
